\section{Colored balanced kernels and their Mellin transform}
\label{sec:colored}

Let
\[
 \Omega=\C\setminus[1,\infty),\qquad a>0,\qquad x>0,\qquad \omega\in\Omega.
\]
The branch of $\Phi(\omega,s,a)$ is determined by its defining series
for $|\omega|<1$ and continuation to $\Omega$. Equivalently, for
$\Re s>0$ it is given by the classical absolutely convergent integral
\cite[25.14.5]{dlmflerch}
\begin{equation}\label{col:lerch-integral}
 \Phi(\omega,s,a)=\frac1{\Gamma(s)}
  \int_0^\infty\frac{t^{s-1}e^{-at}}{1-\omega e^{-t}}\dd t.
\end{equation}
Subtracting Taylor coefficients of its smooth density at zero extends
it to an entire function of $s$ for each $\omega\in\Omega$.

\begin{definition}[Colored balanced derivative]\label{col:def}
For an integer $q\ge0$, put
\begin{align}
 F_\omega(z;x,a)&=x^z\Phi(\omega,1+z,x+a),\label{col:generator}\\
 C_q^\omega(x,a)&=\left.\partial_z^qF_\omega(z;x,a)\right|_{z=0}\notag\\
 &=\sum_{j=0}^q\binom qj\log^{q-j}x\,
       \partial_s^j\Phi(\omega,1,x+a).\label{col:kernel}
\end{align}
Order derivatives are evaluated after analytic continuation in $\omega$.
\end{definition}

There is no order pole at $s=1$ when $\omega\ne1$ in this domain.
Thus adding the uncolored term $\log^{q+1}x/(q+1)$ to
\eqref{col:kernel} would destroy decay rather than improve it.

\subsection{Tail cancellation and uniqueness}

Write
\begin{equation}\label{col:A}
 A_k(\omega,a)=\Phi(\omega,-k,a)
       =(\omega\partial_\omega+a)^k\frac1{1-\omega}.
\end{equation}
These are rational functions of $\omega$ and polynomials of degree $k$
in $a$. Their leading coefficient in $a$ is $1/(1-\omega)$.

\begin{lemma}[Uniform colored expansion]\label{col:tail}
On compact subsets of $\Omega\times(0,\infty)$, and locally uniformly
for $z$ near zero,
\begin{equation}\label{col:F-tail}
 F_\omega(z;x,a)=
  \sum_{k=0}^{K}\frac{(-1)^k A_k(\omega,a)}{k!}
           (1+z)_k x^{-k-1}+O(x^{-K-2}).
\end{equation}
Consequently
\begin{equation}\label{col:C-tail}
 C_q^\omega(x,a)=
  \sum_{k=q}^{K}\frac{(-1)^k A_k(\omega,a)}{k!}
   \left.\partial_z^q(1+z)_k\right|_{z=0}x^{-k-1}
       +O(x^{-K-2}),
\end{equation}
and the leading coefficient is
$(-1)^q A_q(\omega,a)x^{-q-1}$.
At zero, $C_q^\omega(x,a)=O(1+|\log x|^q)$.
\end{lemma}

\begin{proof}
The density $e^{-at}/(1-\omega e^{-t})$ is analytic at $t=0$,
with Taylor coefficients $(-1)^k A_k(\omega,a)/k!$.
Taylor's formula on a small fixed interval and exponential decay
outside that interval give a Laplace expansion in
\eqref{col:lerch-integral}. The $k$th term integrates to
$\Gamma(1+z+k)x^{-1-z-k}$; division by $\Gamma(1+z)$
and multiplication by $x^z$ give \eqref{col:F-tail}.
On the stated compact sets, the denominator stays uniformly away
from zero along $t\ge0$, so the remainders are locally uniform.
Cauchy differentiation in a slightly larger $z$ disc preserves
the inverse-power order of the remainder. Since $(1+z)_k$ has
degree $k$, derivatives of order $q>k$ vanish, and its leading
coefficient gives $\partial_z^q(1+z)_q=q!$.
The assertion at zero follows directly from \eqref{col:kernel}.
\end{proof}

\begin{theorem}[Unique homogeneous balancing]\label{col:unique}
Fix $\omega\in\Omega$ and $q\ge0$. Among combinations
\[
 D(x,a)=\sum_{j=0}^q d_j\log^{q-j}x\,
                  \partial_s^j\Phi(\omega,1,x+a),\qquad d_0=1,
\]
with coefficients independent of $x,a$, the condition
$D(x,a)=O(x^{-q-1})$ for every $a$ in a nonempty open interval contained in $(0,\infty)$
forces $d_j=\binom qj$. These coefficients do satisfy that condition.
\end{theorem}

\begin{proof}
Put $d(t)=\sum_{j=0}^q d_jt^j$ and $L=\log x$.
The coefficient arising from the $k$th density term is a polynomial
in $L$. Its component with $L^{q-r}$ contains
$d^{(r)}(-1)$ times a nonzero scalar and
$\partial_s^r(s)_k|_{s=1}$. For $k=r<q$, the latter derivative
is $r!$. The polynomial $A_r(\omega,a)$ cannot vanish on the
whole prescribed interval. Successively removing $x^{-1},\ldots,x^{-q}$
therefore gives $d^{(r)}(-1)=0$ for $r=0,\ldots,q-1$.
Since $\deg d\le q$ and $d(0)=1$, $d(t)=(1+t)^q$.
The converse is Lemma~\ref{col:tail}. This proves uniqueness only
in the displayed ansatz, not among arbitrary $x$-dependent combinations.
\end{proof}

% Optional insertion after Theorem col:unique. Uses root's notation.
\begin{remark}[The interval in the shift is essential]
\label{col:pointwise-counterexample}
The uniqueness statement can fail if imposed only at one shift.
Set
\[
 a_* = \frac{1+\sqrt2}{2},\qquad
 \omega_*=-3-2\sqrt2\in\Omega.
\]
The rational second moment factors as
\[
 A_2(\omega,a_*)
  =\frac{(3-2\sqrt2)(\omega+3+2\sqrt2)^2}
          {4(1-\omega)^3}.
\]
Consequently $A_2(\omega_*,a_*)=0$ and, by
$A_3=(\omega\partial_\omega+a_*)A_2$,
$A_3(\omega_*,a_*)=0$ as well.
The uniform tail formula therefore gives
\[
 C_2^{\omega_*}(x,a_*)=O(x^{-5}),\qquad
 C_3^{\omega_*}(x,a_*)=O(x^{-5}).
\]
Hence the different homogeneous combination
\[
 D(x)=2C_3^{\omega_*}(x,a_*)
                  -\log x\,C_2^{\omega_*}(x,a_*)
\]
also satisfies $D(x)=O(x^{-4})$.  In the displayed uniqueness
ansatz its coefficient vector is
$\boldsymbol d=(1,4,5,2)$ rather than $(1,3,3,1)$, so it has
the same normalization $d_0=1$.  The fourth moment is nonzero,
$A_4(\omega_*,a_*)=(2-\sqrt2)/64$, and the modified kernel is
not identically the canonical one.  Requiring decay for a
nonempty open interval of positive shifts excludes this
pointwise degeneracy exactly as the proof of the theorem requires.
\end{remark}

\subsection{The exact Mellin identity}

Define the generalized harmonic polynomial $Q_q(u)$ by
\begin{equation}\label{col:Q}
 \frac{\Gamma(u+z)}{\Gamma(u)\Gamma(1+z)}
     =\sum_{q=0}^\infty Q_q(u)\frac{z^q}{q!}.
\end{equation}
With
$H_{u-1}^{(1)}=\psi(u)+\gamma$ and
$H_{u-1}^{(r)}=\zeta(r)-\zeta(r,u)$ for $r\ge2$, this is
\[
 \exp\!\left(\sum_{r=1}^\infty
          \frac{(-1)^{r-1}}r H_{u-1}^{(r)}z^r\right).
\]
In particular,
$Q_0=1$, $Q_1=H^{(1)}$, and
$Q_2=(H^{(1)})^2-H^{(2)}$.

\begin{theorem}[Colored balanced Mellin formula]\label{col:mellin}
For $q\ge0$ and $0<\Re u<q+1$,
\begin{equation}\label{col:mellin-formula}
 \boxed{\displaystyle
 \int_0^\infty x^{u-1}C_q^\omega(x,a)\dd x
  =\frac{\pi}{\sin\pi u}\,Q_q(u)\Phi(\omega,1-u,a).}
\end{equation}
Every apparent singularity at an interior integer is removable.
The integral is absolutely convergent and locally uniform in
$u,\omega,a$ on the stated domains.
\end{theorem}

\begin{proof}
Initially take $0<\Re u<1$ and $z$ sufficiently close to zero.
Insert \eqref{col:lerch-integral} into \eqref{col:generator}.
Absolute Fubini and the Gamma integral give
\begin{align*}
 \int_0^\infty x^{u-1}F_\omega(z;x,a)\dd x
 &=\frac{\Gamma(u+z)}{\Gamma(1+z)}
       \int_0^\infty\frac{t^{-u}e^{-at}}{1-\omega e^{-t}}\dd t\\
 &=\frac{\Gamma(u+z)\Gamma(1-u)}{\Gamma(1+z)}
                  \Phi(\omega,1-u,a).
\end{align*}
The behavior at zero and infinity in Lemma~\ref{col:tail}
justifies locally normal differentiation in $z$.
Euler reflection then proves \eqref{col:mellin-formula}
on the initial strip. The differentiated kernel decays as
$x^{-q-1}$, so its integral is holomorphic on the larger
strip $0<\Re u<q+1$. Analytic continuation establishes the
formula there and makes the interior integer poles removable.
This does not differentiate a divergent integral at an integer.
\end{proof}

\section{Integer colored moments, Stieltjes differences, and primitives}
\label{sec:colored-integer}

For $q\ge m\ge1$, let
\begin{align}
 L_{m,q}
 &:=(-1)^q q!\sum_{\ell=m}^\infty
              \frac{\binom{\ell-1}{m-1}}{\ell^{q+1}}
 \label{col:L-positive}\\
 &=(-1)^{m+q+1}q!
       \sum_{j=0}^{m-1}(-1)^j e_j(1,\tfrac12,\ldots,\tfrac1{m-1})
                         \zeta(q+1-j).\label{col:L-finite}
\end{align}
The first representation is convergent and gives its sign.
The second is a finite zeta expression, with the empty-list convention
$e_0=1$.

\begin{theorem}[All ordinary positive-integer moments]\label{col:integer}
For $q\ge m\ge1$,
\begin{equation}\label{col:integer-formula}
 \boxed{\displaystyle
 \int_0^\infty x^{m-1}C_q^\omega(x,a)\dd x
             =L_{m,q}\Phi(\omega,1-m,a).}
\end{equation}
Thus the complete shift dependence is a polynomial of degree
$m-1$, with rational functions of the color as coefficients.
\end{theorem}

\begin{proof}
At $u=m$, the Gamma quotient in \eqref{col:Q} is
$\prod_{j=1}^{m-1}(1+z/j)$, hence $Q_q(m)=0$ for $q\ge m$.
The limit in \eqref{col:mellin-formula} is
$(-1)^m Q_q'(m)\Phi(\omega,1-m,a)$.
Differentiating the quotient in $u$ gives
\[
 Q_q'(m)=q![z^q]\prod_{j=1}^{m-1}(1+z/j)
                      \bigl(\psi(m+z)-\psi(m)\bigr).
\]
The rational terms in the digamma recurrence become polynomials
of degree at most $m-1$ after multiplication by the finite product.
They do not contribute at degree $q$. Expanding
$\psi(1+z)-\psi(1)=\sum_{r\ge1}(-1)^{r+1}\zeta(r+1)z^r$
proves $(-1)^mQ_q'(m)=L_{m,q}$ in the finite form.
Expanding the zeta values and using
\[
 \prod_{j=1}^{m-1}(1-\ell/j)
       =(-1)^{m-1}\binom{\ell-1}{m-1}
\]
gives \eqref{col:L-positive}.
\end{proof}

There is an independent elementary proof for $|\omega|<1$.
In that disk,
\begin{equation}\label{col:elementary-series}
 C_q^\omega(x,a)=\sum_{\ell=0}^\infty
  \frac{\omega^\ell}{x+\ell+a}
           \log^q\!\frac{x}{x+\ell+a}.
\end{equation}
For $c>0$, the substitution $t=x/(x+c)$ gives
\begin{equation}\label{col:elementary-moment}
 \int_0^\infty\frac{x^{m-1}}{x+c}
                  \log^q\!\frac{x}{x+c}\dd x
 =c^{m-1}\int_0^1\frac{t^{m-1}\log^q t}{(1-t)^m}\dd t
 =c^{m-1}L_{m,q}.
\end{equation}
Absolute summation proves the theorem in the disk, and locally normal
continuation in $\omega$ gives the rest of $\Omega$. At $\omega=-1$,
one must not replace that continuation by the divergent series
$\sum(-1)^\ell(\ell+a)^{m-1}$.

\subsection{Low-order examples and logarithmic insertions}

For every $q\ge1$ the unweighted moment is independent of the shift:
\begin{equation}\label{col:unweighted}
 \boxed{\displaystyle
 \int_0^\infty C_q^\omega(x,a)\dd x
       =\frac{(-1)^q q!\zeta(q+1)}{1-\omega}.}
\end{equation}
For example,
\begin{align}
 \int_0^\infty C_1^\omega(x,a)\dd x
      &=-\frac{\zeta(2)}{1-\omega},\label{col:first}\\
 \int_0^\infty xC_2^\omega(x,a)\dd x
      &=2\bigl(\zeta(2)-\zeta(3)\bigr)
         \left(\frac{a}{1-\omega}
                    +\frac{\omega}{(1-\omega)^2}\right).
 \label{col:second}
\end{align}
All fixed logarithmic insertions are equally explicit. For $q\ge m$
and $r\ge0$,
\begin{equation}\label{col:log-moments}
\begin{split}
 &\int_0^\infty x^{m-1}\log^r x\,C_q^\omega(x,a)\dd x\\
 &\quad=(-1)^m r![t^r]\left\{
   \frac{\pi t}{\sin\pi t}\,
   \frac{Q_q(m+t)}t\,\Phi(\omega,1-m-t,a)\right\}.
\end{split}
\end{equation}
The quotient $Q_q(m+t)/t$ is holomorphic. This is a finite coefficient
formula in Lerch order derivatives, Gamma derivatives, and generalized
harmonic polynomials. It follows by differentiating the normally
convergent Mellin integral, not by termwise integration of its
individual divergent logarithmic summands.

\subsection{A precise generalized Stieltjes dictionary}

Let $\omega^d=1$ with $\omega\ne1$. Introduce the incoming uncolored
kernel \cite{repobalanced}
\begin{equation}\label{col:uncolored}
 C_q(x,a)=\frac{\log^{q+1}x}{q+1}
       +\sum_{j=0}^q(-1)^j\binom qj
                  \gamma_j(x+a)\log^{q-j}x.
\end{equation}
The finite residue-class identity for Lerch's function gives
\begin{equation}\label{col:stieltjes-dictionary}
 \boxed{\displaystyle
 C_q^\omega(x,a)=\frac1d\sum_{r=0}^{d-1}
       \omega^r C_q\!\left(\frac xd,\frac{a+r}{d}\right).}
\end{equation}
Indeed
$\Phi(\omega,s,x+a)=d^{-s}\sum_r\omega^r
\zeta(s,(x+a+r)/d)$; the $1/z$ poles and the
$\log^{q+1}(x/d)$ terms cancel because $\sum_r\omega^r=0$.
This identity translates every preceding colored formula into an
ordinary integral of a finite combination of generalized Stieltjes
constants. It also verifies the normalization against the incoming
finite-filter calculus.

For the alternating color, define
\[
 a_-=\frac{x+a}{2},\qquad a_+=\frac{x+a+1}{2}.
\]
Then, explicitly,
\begin{equation}\label{col:alternating-C1}
 C_1^{-1}(x,a)=\frac12\left[
    \bigl(\psi(a_+)-\psi(a_-)\bigr)\log\frac x2
                 +\gamma_1(a_+)-\gamma_1(a_-)\right].
\end{equation}
Its ordinary integral is $-\pi^2/12$ for every $a>0$.

\subsection{Exceptional moments and Dirichlet beta values}

\begin{theorem}[A zero of the leading colored tail]\label{col:exceptional}
If $m\ge1$ and $\Phi(\omega,1-m,a)=0$, then
\begin{equation}\label{col:exceptional-formula}
 \boxed{\displaystyle
 \int_0^\infty x^{m-1}C_{m-1}^\omega(x,a)\dd x
       =(-1)^{m+1}\partial_s\Phi(\omega,1-m,a).}
\end{equation}
The integral is ordinary and absolutely convergent.
\end{theorem}

\begin{proof}
The zero hypothesis removes the leading $x^{-m}$ kernel term in
Lemma~\ref{col:tail}, so the Mellin integral extends
holomorphically across $u=m$. At $u=m$, $Q_{m-1}(m)=1$ and
$\Phi(\omega,1-u,a)=-(u-m)\partial_s\Phi(\omega,1-m,a)
 +O((u-m)^2)$. The removable limit in
Theorem~\ref{col:mellin} is \eqref{col:exceptional-formula}.
\end{proof}

Let $\beta(s)=\sum_{n\ge0}(-1)^n(2n+1)^{-s}$ denote Dirichlet's beta
function, continued in its order. Since
$\Phi(-1,s,1/2)=2^s\beta(s)$, its negative odd zeros give the following
family:
\begin{corollary}[Alternating centered exceptional moments]\label{col:beta}
For every integer $r\ge1$,
\begin{equation}\label{col:beta-formula}
 \boxed{\displaystyle
 \int_0^\infty x^{2r-1}C_{2r-1}^{-1}(x,\tfrac12)\dd x
   =\frac{(-1)^r(2r-1)!}{\pi^{2r-1}}\beta(2r).}
\end{equation}
In particular, writing $G=\beta(2)$ for Catalan's constant,
\begin{equation}\label{col:catalan}
 \int_0^\infty xC_1^{-1}(x,\tfrac12)\dd x=-\frac G\pi.
\end{equation}
\end{corollary}

\begin{proof}
The completed beta function
$(4/\pi)^{(s+1)/2}\Gamma((s+1)/2)\beta(s)$ is invariant under
$s\mapsto1-s$ \cite{dlmfdirichlet}. Comparing its finite values at
$s=1-2r$ and $s=2r$ yields
\[
 \beta'(1-2r)=
 (-1)^{r-1}2^{2r-1}(2r-1)!\pi^{1-2r}\beta(2r).
\]
Substitute this into Theorem~\ref{col:exceptional}.
The classical functional equation is the input; the ordinary
Stieltjes-difference integral is its consequence in the new kernel.
\end{proof}

\subsection{Anchored shift primitives}

For $a,a_0>0$ the quotient below is regular at $z=0$, because
$\Phi(\omega,0,A)=1/(1-\omega)$ does not depend on $A$:
\begin{equation}\label{col:primitive}
\begin{split}
 \int_{a_0}^a C_q^\omega(x,v)\dd v
  =-q![z^q]\,\frac{x^z}{z}
       \left(\Phi(\omega,z,x+a)-\Phi(\omega,z,x+a_0)\right).
\end{split}
\end{equation}
This follows from
$\partial_A\Phi(\omega,z,A)=-z\Phi(\omega,z+1,A)$.
It gives every integration constant through the explicit anchor.
At order zero and a root of unity, the Hurwitz identity
$\zeta'(0,A)=\log\Gamma(A)-\tfrac12\log(2\pi)$ gives the finite
Gamma expression
\begin{equation}\label{col:gamma-primitive}
 \int_{a_0}^a C_0^\omega(x,v)\dd v
  =-\sum_{r=0}^{d-1}\omega^r
       \log\frac{\Gamma((x+a+r)/d)}
                    {\Gamma((x+a_0+r)/d)}.
\end{equation}
All Gamma arguments are positive, and their logarithms are real
before multiplication by the complex coefficients.

\section{The color-one limit and its exact constant term}
\label{sec:color-limit}

The formulas above are not uniform at $\omega=1$. The precise
normalization can nevertheless be determined, and it reconciles the
colored moment with its uncolored Stieltjes counterpart.
Fix $q\ge m\ge1$, let $\omega=e^{-\lambda}$ with $\lambda>0$, and set
\[
 M_{m,q}(\lambda,a)=\int_0^\infty
                   x^{m-1}C_q^{e^{-\lambda}}(x,a)\dd x.
\]

\begin{theorem}[Complete local color expansion]\label{col:coalescence}
For $0<\lambda<2\pi$,
\begin{equation}\label{col:coalescence-formula}
 \boxed{\displaystyle
 e^{-a\lambda}M_{m,q}(\lambda,a)
 =L_{m,q}\left\{\frac{\Gamma(m)}{\lambda^m}
       +\sum_{k=0}^\infty
             \frac{(-\lambda)^k}{k!}\zeta(1-m-k,a)\right\}.}
\end{equation}
The series is convergent. In the cutoff coordinate $\lambda$,
\begin{align}
 \lim_{\lambda\downarrow0}
 \left(e^{-a\lambda}M_{m,q}
              -L_{m,q}\Gamma(m)\lambda^{-m}\right)
      &=-\frac{L_{m,q}}m B_m(a),\label{col:renorm}\\
 \CT_{\lambda=0}M_{m,q}(\lambda,a)
      &=\frac{L_{m,q}}m\bigl(a^m-B_m(a)\bigr).\label{col:raw-CT}
\end{align}
The first line is exactly the uncolored ordinary moment
$\int_0^\infty x^{m-1}C_q(x,a)\dd x$ from \cite{repobalanced}.
\end{theorem}

\begin{proof}
The classical convergent Lerch expansion \cite[25.14.3\_1]{dlmflerch}
is
\[
 e^{-a\lambda}\Phi(e^{-\lambda},s,a)
    =\Gamma(1-s)\lambda^{s-1}
       +\sum_{k=0}^\infty\zeta(s-k,a)\frac{(-\lambda)^k}{k!}.
\]
Set $s=1-m$ and use Theorem~\ref{col:integer}.
Its constant term is $L_{m,q}\zeta(1-m,a)
=-L_{m,q}B_m(a)/m$.
Multiplying by $e^{a\lambda}$ adds the constant contribution
$L_{m,q}\Gamma(m)a^m/m!=L_{m,q}a^m/m$ from the pole.
This proves both normalizations.
The incoming uncolored coefficient is
$\kappa_{m,q}=-L_{m,q}/m$, so its moment is the first limit.
\end{proof}

For example, $M_{1,1}=-\zeta(2)/(1-e^{-\lambda})$ is independent of
$a$, and its raw constant term is $-\zeta(2)/2$.
The normalized limit in \eqref{col:renorm}, however, is
$\zeta(2)(a-\tfrac12)$. Both formulas are correct in their stated
normalizations. They cannot be identified by discarding the
$e^{-a\lambda}$ factor.

All logarithmic moment versions follow by differentiating the local
Mellin expression in $u$ before putting $u=m$. Derivatives of
$\lambda^{-u}$ produce explicit powers of $\log\lambda$.
Thus the scalar expansion already specifies the complete mechanism
for those contacts. It does not assert that the pointwise limit of
the colored kernels can be integrated across the entire half-line.

